Figure \ref{fig:EquityPlot} shows MUCP and number of administered routine screens split between the vaccinated and unvaccinated subpopulations, as rates per 100k to allow for direct comparison between the two despite thier differing sizes, at 80\% vaccine uptake. Among the unvaccinated subpopulation, trends align closely with population-wise trends as most cancers occur in this group. These trends are mostly attributable to infection with HPV genotypes included in the nonavalent vaccine, with vaccine-types contributing to 12-16 MUCP /100k in the unvaccinated subpopulation in 2040, while non-vaccine types contribute only 0.7-0.8 MUCP /100k in the same year (across all algorithms other than None). There is an increase from 2026 to 2050 in cancers attribtable to vaccine-types over time, however this is driven mostly by the vaccinated population, which experience cancer at a much lower rate than the unvaccinated population (in 2040 for all algorithms which offer screening to vaccinated individuals: 0.05-0.15 MUCP/100k attributable to vaccine-types, and an equal MUCP/100k attribtable to non-vaccine types, with the MUCP/100k attribtable to non-vaccine-types growing to c. 0.2 MUCP/100k by 2050 while the MUCP/100k attributable to vaccine types remains steady). Within the vaccinated subpopulation, cancer prevalence notably increases over time under all algorithms (including continuation of current 5-yearly screening); this is primarily a result of the young demographics of the vaccinated subpopulation compared to the broader population,such that as the population ages there is naturally an increasing rate of cervical cancer, of which most cases occur after the age of 40. [that last exmplanation is also really a discussion point, it would be nice to itnegrate the explanation with the results though I think] Further genotype-specific results in different subpopulations, including plots, are available in the appendices.  


Among vaccinated individuals, trends show markedly different patterns to those in the full population. Across all screening algorithms, MUCP increases from 2026 to 2050, alebit at values an order of magnitude lower than that in unvacccinated individuals. [DISCUSSION: explained by the vaccinated population being much younger overall than the general population, and especially than the unvaccinated population, so while over lifetime there is still a much lower rate of cancer, rates of cancer in the vaccinated will continue to grow beyond 2050 as they age, until an equilibrium is met when vaccinated individuals are equally distributed across the population pyramid.] 


In unvaccinated individuals, we see most cases are attributable to vaccine-genotypes throughout 2026 to 2050, with the number of cases attributable to vaccine types between 10-20 times those attributable to non-vaccine types. Population-wide trends are mostly driven cancers in unvaccinated individuals from vaccine types, so this pattern is very similar to population-wide trends. There are similar rates of cancer in unvaccinated and vaccinated individuals attributable to unvacc-types (but with a decrease over time in unvacc indiviausl while it increases over time in vacc individuals, likely to converge to similar rates over time like due to the different demographics between unvacc and vacc individuals).


Additional to the different overall trends in MUCP in vaccinated indivuduals compared to the general population, our subpopulation results indicate a different relationship between choice of screening algorithm and cancer outcome in this vaccinated subpopulation, which notably appears to depend both on screening intervals for vaccinated and unvaccianted individuals. Expectedly, any complete removal of screening for vaccinated indivduals results in a sharp increase in cancer prevalence among the vaccinated population - however, MUCP among the vaccinated remains below 0.95/100k between 2026 and 2050 even in this case, a factor of 10-20 times less that in the unvaccinated population. Comparing these vaccination-screening-free algorithms, hile 5:None and 5+None:None have similar trajectories, algorithm None (a full removal of screening for all, including unvaccinated individuals) results in the higest cancer prevalence by a signficiant margin. In all three cases, by 2040 cancer cases amongst the vaccinated are mostly attribtable to genotypes not covered in the nonavalent vaccine, however between 2037 and 2050 the rate of MUCP attributable to vaccine-genotypes stabilises around 0.3 for algorithms 5:None and 5+None:None, and around 0.4 for algorithm None. [Discussion: This highlights how screening in one subpopulation affects both HPV and cancer prevelence in populations which mix with that subpopulation. When no unvaccinated nor vaccinated individuals are screened, our model yields observe signficiantly worse outcomes for vaccinated individuals than when only vaccinated indivduals are not screened - as a result of no screening for the unvaccinated, HPV infection is allowed to continue undetected and untreated, allowing for a greater chance of increased prevalence both within the unvacccinated population but also greater chance that interactions between unvaccinated and vaccinated individuals will result in successful transmission. This change in screening only for the unvaccinated subpopulation therefore yields a coupled effect on the vaccinated population. It is notable that this phenomenon emerged organically within HPVsim, particularly as there are only two non-vaccination points of infection control built into our screening model, both of which primary are modelled for treatment of lesions - ablation (for treating of high-grade CIN) and `general cancer treatment' (for treating cervical cancer). Yet even the removal of these treatments in one subpopulation results in a sharp effect in other subpopulations. We note that a limitation of this aspect of our analysis is the assumed uniform mixing with respect to vaccination status; we do not model population clustering by vaccination status. This is justified by the higher rates of all cancers attribtableto any given genotype None compared to 5:None and 5+None:None, and particularly the higher rates of cancers attributable to vaccine-type genotypes, which remain commonly spread between unvaccinated individuals. ] We notice different clustering of algorithm outcomes in the vaccinated subpopulation compared to the unvaccinated subpopulation: all screening algorithms with large screening intervals for vaccinated individuals perform more poorly in terms of cancer prevalence, with algroithms 5:15 and 5+15:15 performing similarly to 10 and 15. Howveer, for algorithms with more modest vaccinated-screening-interval increases - such as 5:7, 5+7, 5:10, and 5+10 - performance is similar to the current baseline algorithm (5). We note again an coupling where unvaccinated screening intervals and outcomes in the vaccinated subpopulation, with algorithm 15 performing signficiantly worse in terms of MUCP between 2040 and 2050 than 10:15 and 10+15:15.  [Discussion: this appears driven by of the time it takes for cervical cancer to develop from a detectible lesion, so there is a pretty steep dropoff in screening quality when increasing a screening interval from 10 to 15 yeasrs [X]. The coupling can be explained as above - even if we are screening vaccinated individuals/cohorts every 15 years, by screening unvaccinated individuals at a shorter interacl we can catch and treat infections in the unvaccianted before they have a chance to spread to as many vaccinated individuals. ] 

[Discussion: In terms of equitability, we therefore see the situations where vaccinated individuals are not screened at all are highly inequitable to the vaccinated, resulting in sharp increases to cancer prevalence in their subgroup which are not captured as clearly when analysing at a population level. Similarly, while algorithms such as 5:15 and 5+15:15 which combine regular screening for the unvaccinated with with 15-year screening intervals for vaccinated individuals/cohorts appear promising, they also result in notably poorer outcomes for vaccinated individuals and may therefore be considered inequitable.]


We also compare probabilities of cervical cancer elimination by 2040, as in Table \ref{tab:EliminationProbs}, when focussing on the vaccinated and unvaccinated subpopulations. Tables are available in the appendices, with key results are decribed here. In the vaccinated subpopulation, elimination probabilities are consistently over 0.97 across all screening uptakes and all algorithms. This includes those algorithms offering no screening to either vaccinated individuals or the full population, however it is lowest for such algorithms; restricting only to those algorithms which offer some screening to the vaccinated, elimination probabilities are uniformly at least 0.99 (2dp). Among unvaccinated individuals, elimination probabilities are signficiantly lower, and depend more heavily on vaccination uptake: for 80\% vaccination uptake, elimination probabilities peak at 0.35 for those algorithms offering the most overall screenings (5, 5:7, 5:10), and elimination probabilities are their lowest at 0.04, for algorithm None.

As all elimination probabiltiies are saturated near 1.00 for the vaccinated subpopulation, it is less useful to directly compare algorithms in terms of these probabilities for vaccinated individuals. However, it is possible to do so for unvaccinated individuals.

As elimination probabilities for unvaccinated individuals are not saturated near 1.00, it is possible to compare algorithms directly in terms of probability of cervical cancer elimination by 2040 among unvaccinated individuals and in terms of robustness to varying vaccination uptake levels. At 60\% uptake, the unvacccinated cancer elimination probability for 5:7 (0.44) and 5:10 (0.44) rounds to 2 decimal places identically to 5 (0.44). [DISCUSSION:a similar point to before about coupling: This appears to be driven by (i) the same interval being in place for these unvaccinated individuals (5 years), and (ii) the screening interval for vaccinated individuals not being so long as to ensure any vaccinated individuals who do happen to be infectious with HPV do get screened and treated before there is a chance for the infection to be spread more widely, and thereby add to the HPV infection, and subsequently cancer, burden in unvaccinated individuals.] The unvaccinated cancer elimination probability at 60\% uptake for 5:15 is a lower 0.4 [indicating that the condition (ii) begins to weaken at a longer screening interval of 15 years, consistent with our other results and literature [X]]. At 60\% uptake, algorithms 5+7:7 (0.41) and 5+10:10 (0.42) have similarly reduced unvaccinated cancer elimination probabilities [, which is explained by the reduction in screening for young unvaccinated individuals under these algorithms.] It is remarkable that this reduction in equity towards younger unvaccinated individuals when comparing 5:7 and 5:10 to 5+7:7 and 5+10:10 is seen already in 2040. The relative disequity further grows from 2040 to 2050 as the vaccinated cohorts constite a greater proportion of the screening-eligible population. For example, at 60\% uptake, unvaccinated cancer elimination has porbability  \(\geq 0.81\) for 5 and 5:7, but 0.76 for 5+7:7, with a similar pattern in 2050 even at higher uptake 80\%, albeit with unvaccinated cancer elimination probabiltiies beginning to saturate near 1.00 at this point. [Therefore, even within this medium-term projection we see an inequity in cohort-wise approaches comapred to vaccination-dependant approaches, and the dispartity between the approaches can be expected to grow in the long-term with a growing proprotion of screening-elibile individuals members of a vaccinated cohort. Therefore, any decision to adopt a cohort-wise approach would warrent future re-evaluation.]

We conclude our equity analysis by noting a somewhat-unintuitive phenomenon that emerges when performing a subgroup analysis splitting by vaccinated and unvacccinated subpopulations, which motivates cautious interpretation when comparing such cancer rates over changing vaccination uptake levels. Taking unvaccinated cancer elimination probabilities by 2040 as an example, and fixing the algorithm at fixed-interval 5-yearly screening ('5'), we observe probabilities 0.44 for 60\% uptake, 0.34 for 80\% uptake, and 0.3 for 90\% uptake: counterintuitively, an increase in vaccination rates causes a decrease in empirical elimination probability. This is explained as follows: although absolute cancer incidence does decrease in the unvaccinated subgroup as vaccination uptake increases, rates per 100k (in 2040) increase. Firstly, cohort-wise herd effect is reached at vaccine-uptake levels under 60\% [https://pmc.ncbi.nlm.nih.gov/articles/PMC12309560/], and as mixing is mostly between individuals in similar age-brackets, we do not expect a signficiantly improved herd effect in unvaccinated individuals as uptake increases from 60\% to 90\%. Further, most cervical cancers occur after the age of 40 [X], meaning that many cancers which do develop between 2026 and 2050 are in the older unvaccinated cohort, and will develop regardless of vaccination uptake in younger individuals (with which there is little mixing). However, the size of the unvaccinated subpopulation itself does shrink notably as vaccination levels vary from 60\% to 90\%, as a result of fewer younger individuals joining the subpopulation (as they take up vaccines). Therefore, when computing rates of cancer incidence in 2040, for example, we therefore observe the numerator (number of cancers developed in 2040) and denominator (size of unvaccinated subpopulation) change disproportionately, and imply increased rates of cancer incidence in the unvaccinated subpopulation as vaccination levels increase, which - although strictly true - is misrepresentative.




